\chapter{Basics}\label{ch:basics}

\section{Cybersecurity}
Through the progression of digitalization, Cybersecurity is of ever increasing relevance to the security of IT-Systems.
An increasing number of devices are connected to the internet, making them part of the cyberspace and therefore potentially vulnerable to cyberattacks.
Almost all sectors of our lives are at risk, when our homes, infrastructure, media devices, industries and governments are connected to the internet.

Cybersecurity aims to protect these devices by achieving the well-known CIA triad, which is an acronym for Confidentiality, Integrity, and Availability. \cite[chapter~3]{Oriyano_2017}
    
\cite[page~20]{Sankaran_Gulasekaran_2021}




%citations
One of the earliest and most commonly used forms of attack used to target access control is a technique known as war driving.
This technique, simply put, is to use a wireless enabled device along with specialized software used to detect or probe wireless networks that have come with in range of the wireless device.
What makes war driving so popular and so effective is the fact that many computer users, both personal and business, have been deploying 802.11 wireless access points for years with little regard for security. 
The deployment of a wireless network that allows the user to roam has taken precedence over taking measures to secure their access points and devices against potential attack.
In response to this deployment of wireless networks that are misconfigured or insecure, we have a category of attackers that engage in this practice known as war driving. 
Those who engage in this activity will build custom rigs built out of a combination of hardware and software, with the intent of cruising public venues to find both wireless access points and devices in some cases. 
These individuals may just be interested in locating and targeting a specific network for their own ends or, in some cases, they may even share this information to an online database or website where any visitor can search for these access points.


When working with any type of network or information one of the fundamental components that needs to be addressed is that of integrity. 
The integrity of data is essential as the receiver of a piece of information needs to have confidence that what they are receiving and relying on is faithful to the intent of the creator and sender of that data. 
If the integrity, and therefore confidence, in data is lost then data becomes essentially worthless.
If you have been in security long enough you have undoubtedly run into the well-known CIA triad, which is shorthand for Confidentiality, Integrity, and Availability.

\cite{Oriyano_2017}